% !TeX program = xelatex
% UTF-8 transcription of all ten source photographs, in handwritten page order.
% Source: Google Drive / Notebook / systemic HTN.
% Original JPEGs are preserved byte-for-byte:
% pic01.jpg = 20261016_225449556.jpg
% pic02.jpg = 20261016_225454935.jpg
% pic03.jpg = 20261016_225500064.jpg
% pic04.jpg = 20261016_225505841.jpg
% pic05.jpg = 20261016_225512385.jpg
% pic06.jpg = 20261016_225517919.jpg
% pic07.jpg = 20261016_225522944.jpg
% pic08.jpg = 20261016_225536662.jpg
% pic09.jpg = 20261016_225542803.jpg
% pic10.jpg = 20261016_225547367.jpg
% Branches in handwritten diagrams are represented by paragraphs and lists.
% Source statements and numerical values are transcribed, not medically revised.
% May be compiled standalone or input into an existing XePersian document.
\ifdefined\settextfont
  \def\SysHTNStandalone{0}
\else
  \def\SysHTNStandalone{1}
  \documentclass[12pt,a4paper]{article}
  \usepackage[margin=2cm]{geometry}
  \usepackage{amsmath,amssymb,array}
  \usepackage{xepersian}
  \IfFontExistsTF{Amiri}{\settextfont{Amiri}}{\settextfont{FreeSerif}}
  \setlatintextfont{FreeSerif}
\fi
\providecommand{\en}[1]{\lr{#1}}
\newcommand{\SysHTNarrow}{\en{$\longleftarrow$}}
\newcommand{\SysHTNunclear}[1]{[خوانش نامطمئن: #1]}
\newcommand{\SysHTNpage}[1]{\clearpage
  \begin{latin}\noindent\textbf{Note book}\hfill Subject: 57\\
  {\small Veor:\enspace\underline{\hspace{1cm}}\quad Month:\enspace\underline{\hspace{1cm}}\quad Date:\enspace\underline{\hspace{1cm}}}\end{latin}
  \noindent صفحهٔ #1\par\medskip}
\newcommand{\SysHTNfoot}{\par\medskip\begin{latin}\noindent\footnotesize BAHARE DANESH ASL\end{latin}}
\ifnum\SysHTNStandalone=1 \begin{document}\fi

\SysHTNpage{۱}
\section*{\en{Systemic Hypertension}}
۱ـ ۶\% بیماران اورژانس با \en{HTN} مراجعه می‌کنند.
\begin{itemize}
\item \en{Severe HTN}.
\item بین ۲۵–۵۰ درصد \SysHTNarrow{} \en{End-organ damage (EOD)}.
\end{itemize}
\subsection*{۲ـ \en{clinic HTN}}
\begin{latin}
\begin{itemize}
\item Prehypertension
\item Stage I HTN
\item Stage II HTN
\end{itemize}
\end{latin}
\subsection*{۳ـ \en{Hypertensive Crisis}}
\begin{latin}
\noindent Acute elevation of blood pressure\\
\fbox{\(SBP\geq180\,\mathrm{mmHg}\) and/or \(DBP\geq120\,\mathrm{mmHg}\)}
\par Forms:
\begin{itemize}
\item Hypertensive Emergency \(\longrightarrow\) HTN \(\uparrow\) + EOD\\
brain, heart, aorta, kidney, eyes.
\item Hypertensive urgency \(\longrightarrow\) HTN -- EOD.
\end{itemize}
\end{latin}
بیمار با هایپرتنسیو ارجنسی را به صورت تهاجمی درمان نمی‌کنیم.
\begin{itemize}
\item درمان ضد فشار خون خوراکی می‌دهیم و پیگیری به صورت \en{Outpatient} می‌کنیم.
\item فشار ظرف چند روز تا چند هفته پایین آورده می‌شود.
\end{itemize}
\subsection*{۴ـ \en{BP Category in adults}}
\begin{latin}
\renewcommand{\arraystretch}{1.3}
\begin{tabular}{llll}
Normal: & \(<120\) & and & \(<80\)\\
Elevated: & 120--129 & and & \(<80\)\\
HTN, Stage I: & 130--139 & or & 80--89\\
HTN, Stage II: & \(\geq140\) & or & \(\geq90\)
\end{tabular}
\end{latin}
\SysHTNfoot

\SysHTNpage{۲}
\section*{۵ـ \en{Hypertensive Emergency}}
\begin{latin}
\subsection*{a) Acute Aortic Dissection}
\noindent Chest pain (CP), back pain.
\begin{itemize}
\item \(>20\,\mathrm{mmHg}\) difference in upper extremities.
\item Evidence of acute EOD:
\begin{itemize}
\item Abnormal CT-angio of chest; Abdomen/Pelvic.
\item Abnormal TE-Echo of aorta.
\end{itemize}
\end{itemize}
\subsection*{b) Acute Pulmonary Edema}
\noindent Shortness of breath.
\par Evidence of acute EOD \(\longrightarrow\) Interstitial Edema on CXR.
\subsection*{c) Acute Myocardial Infarction}
\noindent CP, Nausea/Vomiting, Diaphoresis.
\par Evidence of acute EOD:
\begin{itemize}
\item Abn EKG.
\item Elevated Cardiac biomarkers.
\end{itemize}
\subsection*{d) Acute Coronary Syndrome}
\noindent CP, Nausea/Vomiting, Diaphoresis.
\par Evidence of acute EOD:
\begin{itemize}
\item Clinical diagnosis.
\item Abn EKG.
\item Elevated Cardiac biomarkers.
\end{itemize}
\subsection*{e) Acute Renal Failure}
\noindent Systolic or diastolic Abdominal bruit.
\par Evidence of acute EOD:
\begin{itemize}
\item Elevated Serum Cr.
\item Proteinuria (Abn U/A).
\end{itemize}
\end{latin}
\SysHTNfoot

\SysHTNpage{۳}
\begin{latin}
\subsection*{f) Preeclampsia/Eclampsia}
\begin{itemize}
\item Seizure, shortness of breath, headache.
\item Vision Abnormalities \(\longrightarrow\) blurred vision, flashing lights; scotomata.
\end{itemize}
\end{latin}
\en{scotomata}: (تغییر در بخشی از بینایی \SysHTNunclear{عبارت فارسی کنار \en{scotomata}}).
\begin{latin}
\noindent Evidence of acute EOD:
\begin{itemize}
\item low Plt.
\item Renal insufficiency.
\item Elevated Liver enzyme levels.
\item Pulmonary Edema.
\item Proteinuria (No longer required for preeclampsia).
\end{itemize}
\subsection*{g) Hypertensive retinopathy}
\noindent Blurred Vision.
\par Evidence of acute EOD:
\begin{itemize}
\item Retinal Hemorrhage.
\item Cotton-wool spots.
\item Hard exudates.
\item Sausage-shaped veins.
\end{itemize}
\subsection*{h) Hypertensive encephalopathy}
\begin{itemize}
\item Altered Mental status.
\item Nausea/Vomiting, Headache.
\end{itemize}
\noindent E. of acute EOD:
\begin{itemize}
\item Papilledema or arteriolar hemorrhage or exudates on funduscopy.
\item Cerebral edema on MRI.
\end{itemize}
\end{latin}
\SysHTNfoot

\SysHTNpage{۴}
\begin{latin}
\subsection*{i) SAH}
\noindent Headache, focal neurologic deficits.
\par Evidence of acute EOD:
\begin{itemize}
\item Abn brain CT.
\item RBC on LP.
\end{itemize}
\subsection*{j) ICH}
\noindent Headache, new neurologic deficits.
\par Evidence of acute EOD \(\longrightarrow\) Abn brain CT.
\subsection*{k) Acute Ischemic Stroke}
\noindent New neurologic deficit.
\par Evidence of acute EOD \(\longrightarrow\) Abn brain CT or MRI.
\subsection*{L) Acute perioperative HTN}
\noindent Bleeding unresponsive to direct pressure.
\par E. of acute EOD:
\begin{itemize}
\item Clinical diagnosis.
\item Manifestations of other emergencies.
\end{itemize}
\subsection*{M) Sympathetic Crisis}
\begin{itemize}
\item EOD may not be measurable.
\item Anxiety, palpitation, tachycardia, diaphoresis.
\end{itemize}
\noindent Evidence of acute EOD:
\begin{itemize}
\item Clinical diagnosis in the setting of sympathomimetic drug use \(\longrightarrow\) Cocaine, Amphetamines.
\item 24 h urine assay for catecholamines and metanephrine or plasma fractionated metanephrines \(\longrightarrow\) pheochromocytoma.
\end{itemize}
\end{latin}
\SysHTNfoot

\SysHTNpage{۵}
۶ـ فشار هر دو اندام فوقانی بیمار را در حالتی که بیمار استراحت می‌کند اندازه بگیرید.

۷ـ قبل از شروع درمان ضد فشار خون چندین نوبت \en{BP} را اندازه‌گیری کنید.

۸ـ تفاوت فشار خون بین اندام فوقانی و تحتانی:
\begin{itemize}
\item دایسکشن آئورت.
\item کوآرکتاسیون.
\item بیماری عروق محیطی (\en{PVD}).
\item برخی بیماری‌های یونیلترال عصبی و اسکلتی عضلانی.
\end{itemize}
۹ـ تفاوت فشار خون بین دو اندام فوقانی:
\begin{itemize}
\item ممکن است در افراد مسن طبیعی باشد.
\item اختلاف بیش از \en{10--20 mmHg} بین دو اندام فوقانی معنی‌دار است.
\end{itemize}
۱۰ـ اگر فشار در دو اندام فوقانی متفاوت بود، فشار بالاتر را درمان می‌کنیم.
\begin{itemize}
\item اندازه‌گیری‌های بعدی را در اندام با فشار بالاتر انجام می‌دهیم.
\end{itemize}
۱۱ـ فشار مچ را اندازه نمی‌گیریم چون پایین‌تر نشان می‌دهد.

۱۲ـ اینکه بیمار تاکنون سابقهٔ \en{HTN} نداشته احتمال اورژانس فشار خون را رد نمی‌کند.
\begin{itemize}
\item ۱۶\% هیچ‌گونه سابقهٔ \en{HTN} ندارند.
\end{itemize}
۱۳ـ ممکن است بدون درمان یک کاهش کم (در حد \en{12 mmHg}) در \en{BP} بیمار مراجعه‌کننده با فشار خون رخ دهد.
\SysHTNfoot

\SysHTNpage{۶}
\subsection*{۱۴ـ \en{CP and severe HTN}}
\begin{latin}
\noindent DDX of aortic dissection from acute coronary syndromes is imperative.
\begin{itemize}
\item BP control differs.
\item Anticoagulation is catastrophic in aortic dissection.
\end{itemize}
\noindent Acute aortic dissection:
\begin{itemize}
\item Abrupt, sudden onset of pain.
\item Pain is usually in the chest \(\longrightarrow\) Tearing or ripping.
\item Radiating to the interscapular region.
\end{itemize}
\end{latin}
\subsection*{۱۵ـ \en{Acute neurologic symptoms and severe HTN}}
\begin{latin}
\noindent Ischemic or hemorrhagic stroke \(\longrightarrow\) headache, focal neurologic deficit.
\par Hypertensive encephalopathy:
\begin{itemize}
\item A clinical diagnosis made after excluding focal ischemia or bleeding.
\item Altered mental status, headache, Nausea/Vomiting, Seizure, Visual disturbances.
\item Most patients will have papilledema.
\end{itemize}
\noindent $\star$ When MRI findings demonstrate reversible edema that is predominantly occipital, the Posterior Reversible Encephalopathy Syndrome exists \(\longrightarrow\) Poor prognosis.
\end{latin}
\subsection*{۱۶ـ \en{ARF, Peripheral edema and severe hypertension}}
\begin{latin}
\noindent New-onset renal failure with:
\begin{itemize}
\item peripheral edema, oliguria, loss of appetite, N/V, confusion.
\item Elevated Serum Cr; Urinary sediment is abnormal.
\end{itemize}
\end{latin}
\SysHTNfoot

\SysHTNpage{۷}
\subsection*{۱۷ـ \en{Preeclampsia}}
\begin{latin}
\noindent Elevated BP \(\longrightarrow SBP\geq140\) or \(DBP\geq90\).
\begin{itemize}
\item On 2 occasions at least 4 h apart.
\item In pregnant patients beyond the 20th week.
\end{itemize}
\noindent Eclampsia:
\begin{itemize}
\item New-onset grand mal seizures.
\item In the absence of other neurologic conditions.
\end{itemize}
\end{latin}
\subsection*{۱۸ـ \en{Sympathetic Crisis and severe HTN}}
\begin{latin}
\subsubsection*{a) Acute catecholaminergic syndrome}
\begin{itemize}
\item Abrupt discontinuation of oral or TD Clonidine.
\item Potentiated by concomitant \(\beta\)-blocker therapy due to unopposed \(\alpha\)-mediated vasoconstriction.
\end{itemize}
\subsubsection*{b) Pheochromocytoma}
\begin{itemize}
\item Life-threatening HTN.
\item Asymptomatic periods, episodic headache, elevated BP, tachycardia, diaphoresis.
\end{itemize}
\subsubsection*{c) Sympathomimetic drugs}
\begin{itemize}
\item Cocaine, Amphetamines, phencyclidine hydrochloride, lysergic acid diethylamide.
\item Tachycardia, diaphoresis, chest pain, mental status changes.
\item MAO inhibitors + consume tyramine-containing foods.
\end{itemize}
\end{latin}
\SysHTNfoot

\SysHTNpage{۸}
\begin{latin}
\subsection*{d) Autonomic dysfunction}
\begin{itemize}
\item Due to spinal cord or severe head injury.
\item Spina bifida.
\end{itemize}
\noindent $\star$ Often the associated BP measurements are marginally elevated in this condition \(\longrightarrow\) Near normal value.
\end{latin}
\subsection*{۱۹ـ بیماران بدون علامت با \en{Severe HTN}}
\begin{latin}
\noindent Basic metabolic panel, EKG, chest radiograph, U/A.
\end{latin}
بررسی‌ها در کمتر از ۶\% آن‌ها یافته‌ای دارد.

$\star$ در فشار خون ما بیشتر به علائم بیمار باید توجه داشته باشیم.

۲۰ـ بیمار با اورژانس \en{HTN} باید در \en{Critical Care Setting} بستری شود.

۲۱ـ به طور کلی فشار سیستولیک را در ساعت اول بیشتر از ۲۵\% کاهش ندهید.
\begin{itemize}
\item ظرف ۲–۶ ساعت بعد: کمتر از \en{160/100} برسانید.
\item ظرف ۲۴–۴۸ ساعت بعد: \en{N/L} برسانید.
\end{itemize}
۲۲ـ در کدام بیماری‌ها با سرعت بیشتری فشار را پایین می‌آوریم؟
\begin{latin}
\begin{itemize}
\item Acute aortic dissection.
\item Eclampsia or severe preeclampsia.
\item Pheochromocytoma crisis.
\item Acute ICH.
\end{itemize}
\end{latin}
\SysHTNfoot

\SysHTNpage{۹}
\subsection*{۲۳ـ درمان دایسکشن آئورت}
هدف درمان:
\begin{itemize}
\item \en{SBP: 120--100}.
\item \en{$HR\leq60$}.
\item طی یک ساعت اول.
\end{itemize}
برای کنترل درد اوپیوئید می‌دهیم \SysHTNarrow{} تون سمپاتیک را کاهش می‌دهد.
\subsection*{۲۴ـ درمان ادم ریهٔ ناشی از \en{HTN}}
باید پاتولوژی زمینه‌ای را درمان کنیم.
\begin{latin}
\noindent Most patients have existing poorly controlled HTN with cardiac remodeling and LVH, stiffness and diastolic disfunction.
\par With acute rise in BP:
\begin{itemize}
\item Increase in afterload.
\item Decrease in Venous capacitance.
\end{itemize}
\noindent Leads to fluid shift from splanchnic and peripheral vascular beds \(\longrightarrow\) into the pulmonary circulation.
\end{latin}
کاهش افترلود مؤثرتر از دیورز است.

درمان اصلی وازودیلاتورها هستند.
\begin{itemize}
\item نیترات \SysHTNarrow{} \en{IV} و زیرزبانی، تاپیکال.
\end{itemize}
$\star$ در صورتی که خواستیم دیورتیک بدهیم باید در نظر بگیریم که لوپ دیورتیک‌ها در کنار نیترات عملکرد کلیه را مختل می‌کنند.
\begin{latin}
\noindent In patient with S. disfunction \(\longrightarrow\) IV Nicardipine, Clevidipine.
\end{latin}
\SysHTNarrow{} افزایش حجم ضربه‌ای و جریان خون کرونری.
\SysHTNfoot

\SysHTNpage{۱۰}
\subsection*{۲۵ـ \en{Treatment of acute myocardial infarction}}
\begin{latin}
\noindent Sublingual or IV nitrates.
\end{latin}
\begin{itemize}
\item بتابلاکر وریدی فقط در بیمارانی که \en{Severe HTN} دارند.
\item بتابلاکر خوراکی \SysHTNarrow{} \en{NSTEMI, STEMI}.
\end{itemize}
\subsection*{۲۶ـ درمان بحران سمپاتیک حاد}
\subsubsection*{\en{a)} مصرف کوکائین یا آمفتامین}
\begin{itemize}
\item بنزودیازپین وریدی.
\item اگر بنزودیازپین جواب نداد \SysHTNarrow{} نیتروگلیسیرین یا فنتولامین.
\item \en{CCB} خط سوم درمان.
\item بتابلاکر نمی‌دهیم چون با اثر آلفابلاکری منجر به وازودیلاته شدن عروق کرونری می‌شود.
\item اگر خواستیم بدهیم \SysHTNarrow{} لابتالول + وازودیلاتور.
\end{itemize}
\subsubsection*{\en{b)} فئوکروموسایتوم}
\begin{itemize}
\item خط اول فنتولامین وریدی است.
\item اگر دسترسی \en{IV} نبود \en{IM} می‌دهیم.
\item خط دوم \SysHTNarrow{} نیکاردیپین یا کلویدیپین.
\end{itemize}
\subsubsection*{\en{c)} مسمومیت با \en{MAOI}}
\begin{itemize}
\item بنزودیازپین وریدی.
\item اگر کافی نبود:
\begin{itemize}
\item فنتولامین.
\item نیتروگلیسیرین.
\item نیتروپروساید.
\end{itemize}
\end{itemize}
\SysHTNfoot

\ifnum\SysHTNStandalone=1
  \let\SysHTNEnd\enddocument
\else
  \let\SysHTNEnd\relax
\fi
\SysHTNEnd
